\documentclass[11pt,oneside]{article}
\usepackage[utf8]{inputenc}
\usepackage{import}
\usepackage[a4paper, bottom=2cm,left=2.5cm,right=2.5cm,top=2.5cm]{geometry}
\usepackage[german]{babel}
\usepackage{framed}
\usepackage{tikz,graphicx}
\usepackage{multicol}
\usepackage{multirow}
\usepackage{xcolor, colortbl}
\usepackage[german,ruled,vlined,linesnumbered,algosection]{algorithm2e}
\usepackage[labelfont=sf]{caption}
\usepackage{enumerate}
\usepackage{enumitem}
\usepackage{float} % Allows putting an [H] in \begin{figure} to specify the exact location of the figure
\usepackage{wrapfig} % Allows in-line images such as the example fish picture
\usepackage[position=top]{subfig}
\RequirePackage{amssymb, amsfonts, amsmath, latexsym, verbatim}
\usepackage{amsthm}
\usepackage{mathtools}
\usepackage{lipsum} % Used for inserting dummy 'Lorem ipsum' text into the template
\usepackage{lmodern}
\usepackage{lscape}
\usepackage{tabularx}
\usepackage{systeme}
\usepackage{hhline}
\usepackage{subfiles}
\usepackage{stackengine}
\usepackage{MnSymbol}
\usepackage[onehalfspacing]{setspace}
\usepackage{tcolorbox}
\usepackage[colorinlistoftodos]{todonotes}
\usepackage{bbm}
\usepackage{fancyhdr}
\usepackage{url}
%\usepackage{tkz-berge}
\usepackage{xparse}
\usepackage{sectsty}
\usepackage{hyperref}
\usepackage{array}
\usepackage{ragged2e}

% definiert die Sprache
\selectlanguage{German}
%\setlocalcaption{german}{figure}{Figur}
%\setlocalcaption{german}{table}{Tabelle}


\tcbuselibrary{breakable,skins}
\usetikzlibrary{mindmap,trees}
\newlength\tindent
\setlength{\tindent}{\parindent}
\setlength{\parindent}{0pt}
\renewcommand{\indent}{\hspace*{\tindent}}
\makeatletter
\let\org@subfile\subfile
\renewcommand*{\subfile}[1]{%
  \filename@parse{#1}% LaTeX's file name parser
  \expandafter
  \graphicspath\expandafter{\expandafter{\filename@area}}%
  \org@subfile{#1}%
}
\makeatother

\makeatletter
\def\ifEmpty#1{\def\@temp{#1}\ifx\@temp\@empty}
\makeatother

\sectionfont{\underline}
\usetikzlibrary{arrows,%
                petri,%
                topaths, graphs}%
\newcommand{\re}[1]{\scriptscriptstyle \mathrm{#1}}

% =================== BEISPIEL ====================
\theoremstyle{definition}
\newtheorem{beispiel}{Beispiel}[section]
\renewcommand{\thebeispiel}{\arabic{beispiel}}

% =================== AUFGABE =====================
\newcounter{aufg}[section] % Zähler für die Umgebungen

\newenvironment{aufgabe}[1][]{%
    \refstepcounter{aufg}% Inkrementiere den Zähler
    \begin{tcolorbox}[%
        enhanced,
        overlay unbroken and first={\mytitleaufg[#1]}, % Titelleiste mit Icons wird nicht über Seite gebrochen
        colframe=black,
        boxrule=0pt,
        breakable,
        top=15pt,
        before=\vskip18pt,
    ]
}{%
    \end{tcolorbox}
}

\newcommand{\mytitleaufg}[1][]{% Definiert die Titelleiste und Icons
    \node[fill=black!20,
        rounded corners,
        draw=black,
        text=black,
        line width=1pt,
        inner sep=4pt,
        anchor=west,
        xshift=12pt]
        at (frame.north west){\bfseries Aufgabe \thesection.\theaufg.\ifEmpty{#1}\else\emph{ (#1)}\fi}; % Titel, letzter Teil mit ifempty testet ob #1 leer und macht Klammern und Titel, falls Titelargument übergeben wird bzw. falls #1 leer erscheint nichts - auch keine Klammern

    \node[% Icon
        rounded corners,
        text=black,
        fill=black,
        line width=0pt,
        inner sep=3pt,
        anchor=west,
        xshift=5pt,
        % yshift=-35pt
        ]
        at (frame.north east){
\includegraphics[width=24pt]{../../../Header/Icons/aufgabeicon.png}};
}

% =================== DEFINITION ==================
\newcounter{def}[section] % Zähler für die Umgebungen

\newenvironment{definition}[1][]{%
    \refstepcounter{def}% Inkrementiere den Zähler
    \begin{tcolorbox}[%
        enhanced,
        drop shadow southeast,
        overlay unbroken and first={\mytitledef[#1]}, % Titelleiste mit Icons wird nicht über Seite gebrochen
        colback=blue!5!white,
        colframe=blue!75!black,
        boxrule=2pt,
        breakable,
        top=15pt,
        before=\vskip18pt,
    ]
}{%
    \end{tcolorbox}
}

\newcommand{\mytitledef}[1][]{% Definiert die Titelleiste und Icons
    \node[fill=blue!40!white,
        rounded corners,
        draw=black,
        text=black,
        line width=1pt,
        inner sep=4pt,
        anchor=west,
        xshift=12pt]
        at (frame.north west){\bfseries Definition \thesection.\thedef.\ifEmpty{#1}\else\emph{ (#1)}\fi}; % Titel, letzter Teil mit ifempty testet ob #1 leer und macht Klammern und Titel, falls Titelargument übergeben wird bzw. falls #1 leer erscheint nichts - auch keine Klammern

    \node[% Icon
        rounded corners,
        text=black,
        fill=black,
        line width=0pt,
        inner sep=3pt,
        anchor=west,
        xshift=5pt,
        % yshift=-35pt
        ]
        at (frame.north east){
\includegraphics[width=24pt]{../../../Header/Icons/definitionicon.png}};
}

% =================== NOTATION ====================
\newcounter{nota}[section] % Zähler für die Umgebungen

\newenvironment{notation}[1][]{%
    \refstepcounter{nota}% Inkrementiere den Zähler
    \begin{tcolorbox}[%
        enhanced,
        drop shadow southeast,
        colback=yellow!5!white,
        colframe=yellow!75!black,
        overlay unbroken and first={\mytitlenota[#1]},   % Titelleiste mit Icons wird nicht über Seite gebrochen
        boxrule=2pt,
        breakable,
        top=15pt,
        before=\vskip18pt,
    ]
}{%
    \end{tcolorbox}
}

\newcommand{\mytitlenota}[1][]{% Definiert die Titelleiste und Icons
    \node[fill=yellow!40!white,,
        rounded corners,
        draw=black,
        text=black,
        line width=1pt,
        inner sep=4pt,
        anchor=west,
        xshift=12pt]
        at (frame.north west){\bfseries Notation \thesection.\thenota.\ifEmpty{#1}\else\emph{ (#1)}\fi}; % Titel, letzter Teil mit ifempty testet ob #1 leer und macht Klammern und Titel, falls Titelargument übergeben wird bzw. falls #1 leer erscheint nichts - auch keine Klammern

    \node[% Icon
        rounded corners,
        text=black,
        fill=black,
        line width=0pt,
        inner sep=3pt,
        anchor=west,
        xshift=5pt,
        % yshift=-35pt
        ]
        at (frame.north east){
\includegraphics[width=24pt]{../../../Header/Icons/definitionicon.png}};
}

% =================== SATZ ========================
\newcounter{satz}[section]  % creates a counter and resets it after each chapter

\newenvironment{satz}[1][]{   % Hauptbox mit Satz drin
    \refstepcounter{satz}
    \begin{tcolorbox}[%
        enhanced,
        colback=red!5!white,
        fuzzy halo=2mm with gray,
        colframe=red!75!black,
        overlay unbroken and first={\mytitlesatz[#1]},   % Titelleiste mit Icons wird nicht über Seite gebrochen
        boxrule=2pt,
        breakable,
        top=15pt,
        before=\vskip18pt,
    ]
}{%
    \end{tcolorbox}
}

\newcommand{\mytitlesatz}[1][]{                     % Macht Titelbox und Icons
   \node[fill=red!90!white,
      rounded corners,
      draw=black,,
      text=black,
      line width=1pt,
      inner sep=4pt,
      anchor=west,
      xshift=12pt]
   at (frame.north west){\bfseries Satz \thesection.\thesatz.\ifEmpty{#1}\else\emph{ (#1)}\fi};  % Titel, letzter Teil mit ifx testet ob #1 leer und macht Klammern und Titel, falls Titelargument übergeben wird bzw.  falls #1 leer erscheint nichts - auch keine Klammern

 \node[ % Icon
      rounded corners,
      text=black,
      fill=black,
      line width=0pt,
      inner sep=3pt,
      anchor=west,
      xshift=5pt,
%     yshift=-35pt
]
   at (frame.north east){ 
\includegraphics[width=24pt]{../../../Header/Icons/satzicon.png} };
}

% =================== BEWEIS ======================
\newenvironment{beweis}[1][]{
    \begin{tcolorbox}[
        enhanced,
        overlay unbroken and first={\mybeweis[#1]},     % Titelleiste mit Icons wird nicht über Seite gebrochen
        colback=white,
        colframe=black,
        boxrule=0pt,
        breakable,
        top=15pt,
        before=\vskip18pt,
    ]
}{
        \end{tcolorbox}
}

\newcommand{\mybeweis}[1][]{                     % Macht Titelbox und Icons
   \node[fill=white,
      rounded corners,
      draw=black,
      text=black,
      line width=1pt,
      inner sep=4pt,
      anchor=west,
      xshift=12pt]
   at (frame.north west){\bfseries Beweis \ifEmpty{#1}\else\emph{ (#1)}\fi};  % Titel, letzter Teil mit ifx testet ob #1 leer und macht Klammern und Titel, falls Titelargument übergeben wird bzw.  falls #1 leer erscheint nichts - auch keine Klammern
}

% =================== ALGORITHM ===================
\newcounter{algo}[section]  % creates a counter and resets it after each chapter

\newenvironment{algo}[1][]{   % Hauptbox mit Satz drin
    \refstepcounter{algo}
    \begin{tcolorbox}[%
        enhanced,
        drop shadow southeast,
        colback=yellow!5!white,
        colframe=yellow!75!black,
        overlay unbroken and first={\algotitle[#1]},   % Titelleiste mit Icons wird nicht über Seite gebrochen
        boxrule=2pt,
        breakable,
        top=15pt,
        before=\vskip18pt,
    ]
}{%
    \end{tcolorbox}
}

\newcommand{\algotitle}[1][]{                     % Macht Titelbox und Icons
   \node[fill=yellow!40!white,
      rounded corners,
      draw=black,
      text=black,
      line width=1pt,
      inner sep=4pt,
      anchor=west,
      xshift=12pt]
   at (frame.north west){\bfseries \ifEmpty{#1}{Ablauf}\else{#1}\fi};  % Titel, letzter Teil mit ifx testet ob #1 leer und macht Klammern und Titel, falls Titelargument übergeben wird bzw.  falls #1 leer erscheint nichts - auch keine Klammern

 \node[ % Icon
      rounded corners,
      text=black,
      fill=black,
      line width=0pt,
      inner sep=3pt,
      anchor=west,
      xshift=5pt
      % yshift=-35pt
      ]
   at (frame.north east){ 
\includegraphics[width=24pt]{../../../Header/Icons/algoicon.png} };
}

% ================ Nummerierte Multicol ================
\newenvironment{enumcols}[1][]{
    \begin{enumerate}[label=\alph*)]\begin{multicols}{#1}
}{
    \end{multicols}\end{enumerate}
}


%%
% Header and Footer
%%
\linespread{1.2} % Line spacing
\fancypagestyle{plain}{
   \pagestyle{fancy}
\lhead{}
\chead{}
\rhead{}
\lfoot{\Autor}
\cfoot{\Skriptname}
\rfoot{\thepage}
\renewcommand{\headrulewidth}{0pt}
\renewcommand{\footrulewidth}{0.4pt}
}

\pagestyle{fancy}
\fancyhf{}
\rhead{\rightmark}
\lfoot{\leftmark}
\rfoot{\thepage}
\renewcommand{\headrulewidth}{0.2pt}
%\setlength\parindent{0pt} % Uncomment to remove all indentation from paragraphs

\usetikzlibrary{shapes,arrows}
\tikzstyle{decision} = [diamond, draw, fill=blue!20,
    text width=6em, text badly centered, node distance=4cm, inner sep=0pt]
\tikzstyle{block} = [rectangle, draw, fill=blue!20,
    text width=8em, text centered, rounded corners, minimum height=4em]
\tikzstyle{line} = [draw, -latex']
\tikzstyle{cloud} = [draw, ellipse,fill=red!20, node distance=2cm,
    minimum height=2em]

%%
% Math Arrows
%%
\newcommand{\same}{\ensuremath{\Longleftrightarrow}}
\renewcommand{\to}{\ensuremath{\longrightarrow}}
\newcommand{\qsame}{\quad\same\quad}
\newcommand{\dsame}{\ensuremath{:\Leftrightarrow}}
\newcommand{\qdsame}{\quad\dsame\quad}
\newcommand{\ldsame}{\ensuremath{:\Longleftrightarrow}}
\newcommand{\samedef}{\ensuremath{\us {\text{def}} \Leftrightarrow}}
\newcommand{\qsamedef}{\quad\samedef\quad}
\newcommand{\lsamedef}{\ensuremath{\us {\text{def}} \Longleftrightarrow}}
\newcommand{\qlsamedef}{\quad\lsamedef\quad}

%%
% Mengendefinitionen
%%
\newcommand{\M}{\ensuremath{\mathcal{M}}}
\newcommand{\N}{\ensuremath{\mathbb{N}}}
\newcommand{\Z}{\ensuremath{\mathbb{Z}}}
\newcommand{\Q}{\ensuremath{\mathbb{Q}}}
\newcommand{\I}{\ensuremath{\mathbb{I}}}
\newcommand{\R}{\ensuremath{\mathbb{R}}}
\newcommand{\C}{\ensuremath{\mathbb{C}}}
\renewcommand{\S}{\ensuremath{\mathbb{S}}}
\newcommand{\Prim}{\ensuremath{\mathbb{P}}}
\newcommand{\0}{\ensuremath{\emptyset}}
\renewcommand{\Re}{\ensuremath{\operatorname{Re}}}
\renewcommand{\Im}{\ensuremath{\operatorname{Im}}}
\newcommand{\strich}{\ensuremath{\big\vert}}

%%
% Logisches Und und Oder
%%
\newcommand{\sand}{\ensuremath{\; \land \;}}
\newcommand{\sor}{\ensuremath{\; \lor \;}}
\renewcommand{\*}{\ensuremath{\cdot}}
\newcommand{\oder}{\vee}
\newcommand{\n}{\neg}
\newcommand{\und}{\wedge}

%%
% nicht Element von
%%
\renewcommand{\nin}{\not\in}

%%
% Bei Integral, Summe,... die Grenzen über bzw.
% unter das Zeichen
%%
\renewcommand{\d}[1]{\ensuremath{\operatorname{d}\!{#1}}}
\newcommand{\Sum}{\sum\limits}
\newcommand{\Prod}{\prod\limits}
\newcommand{\Min}{\min\limits}
\newcommand{\Max}{\max\limits}
\newcommand{\Lim}{\lim\limits}
\newcommand{\Inf}{\inf\limits}
\newcommand{\Sup}{\sup\limits}
\newcommand{\Liminf}{\liminf\limits}
\newcommand{\Limsup}{\limsup\limits}
%\renewcommand{\Cup}{\bigcup\limits}
%\renewcommand{\Cap}{\bigcap\limits}
\newcommand{\Otimes}{\bigotimes\limits}
\newcommand{\Oplus}{\bigoplus\limits}


%%
% Arcus Cotangens, Logarithmus
%%
\let\oldsin\sin
\renewcommand{\sin}[1]{\oldsin{\left(#1\right)}}
\let\oldcos\cos
\renewcommand{\cos}[1]{\oldcos{\left(#1\right)}}
\let\oldtan\tan
\renewcommand{\tan}[1]{\oldtan{\left(#1\right)}}
\let\oldarcsin\arcsin
\renewcommand{\arcsin}[1]{\oldarcsin{\left(#1\right)}}
\let\oldarccos\arccos
\renewcommand{\arccos}[1]{\oldarccos{\left(#1\right)}}
\let\oldarctan\arctan
\renewcommand{\arctan}[1]{\oldarctan{\left(#1\right)}}
\let\oldln\ln
\renewcommand{\ln}[1]{\oldln{\left(#1\right)}}
\newcommand{\Log}[2]{\ensuremath{\operatorname{log}_{#1}\left(#2\right)}}
\newcommand{\e}{\operatorname{e}}
\renewcommand{\exp}[1]{\ensuremath{\operatorname{\e}^{#1}}}


%%
% Griechische Buchstaben
%%
\renewcommand{\epsilon}{\ensuremath{\varepsilon}}
\renewcommand{\phi}{\varphi}
\renewcommand{\l}{\lambda}
\renewcommand{\t}{\tau}


%%
% Bezeichnungen
%%
\newcommand{\Rgn}{\R_{>0}} % R grösser Null
\newcommand{\Rad}{\operatorname{Rad}}   % Radikal
\newcommand{\kgV}{\operatorname{kgV}}   % Kleinster gemeinsame Vielfache
\newcommand{\ggT}{\operatorname{ggT}}   % Grösster gemeinsame Teiler
\newcommand{\winkel}{\sphericalangle}          % Winkelsymbol
\DeclarePairedDelimiter\abs{\lvert}{\rvert} % besserer Betrag
\makeatletter
\let\oldabs\abs
\def\abs{\@ifstar{\oldabs}{\oldabs*}}
\makeatother

%%
% Vektorgeometrie
%%
\renewcommand{\v}[1]{\vv{#1}}       % besserer Pfeil über Buchstaben
\renewcommand{\vec}[2]{\begingroup\renewcommand{\arraystretch}{1}\left(\begin{array}{c} #1 \\ #2 \end{array}\right)\endgroup}   % two dimensional vector
\renewcommand{\Vec}[3]{\begingroup\renewcommand{\arraystretch}{1}\left(\begin{array}{c} #1 \\ #2 \\ #3 \end{array}\right)\endgroup}      % three dimensional vector

%%
% Text
%%
\newcommand*{\Ueq}[1]{\mathrel{\mathop{=}\limits_{\text{#1}}}}
\newcommand*{\Oeq}[1]{\mathrel{\stackrel{\text{#1}}{=}}}
\newcommand{\utext}[2]{\underbrace{#1}_{\substack{#2}}}
\newcommand{\otext}[2]{$\overbrace{\textrm{#1}}^{\textrm{#2}}$}
\newcommand{\lineortext}[1]{\hspace{0,1cm}\underline{\hspace{0.1cm}\hphantom{#1}\hspace{0.1cm}}\hspace{0,1cm}} % Lückentext erstellen
\newcommand{\gap}[1]{\rule{#1 cm}{1pt}}


% add command for different dates in Swiss fashion
\newcommand{\monthwordlong}[1]{\ifcase#1 \or Januar\or Februar\or März\or April\or Mai\or Juni\or Juli\or August\or September\or Oktober\or November\or Dezember\fi}
\newcommand{\monthwordshort}[1]{\ifcase#1\or Jan\or Feb\or Mar\or Apr\or Mai\or Jun\or Jul\or Aug\or Sept\or Okt\or Nov\or Dez\fi}
\newcommand{\datelong}{\number\day.\:\monthwordlong{\the\month}\:\number\year}
\newcommand{\dateshort}{\number\day.\:\monthwordshort{\the\month}\:\number\year}

% define color of links inside a text
\definecolor{linkblue}{RGB}{0, 0, 80}
\definecolor{plotblue}{RGB}{100, 100, 255}
\definecolor{myblue}{RGB}{8,164,214}
\definecolor{myred}{RGB}{143,42,41}
\definecolor{forestgreen}{RGB}{34,139,34}
\definecolor{highdive}{RGB}{10,169,194}
\hypersetup{
	colorlinks=true,
	linkcolor=linkblue,
	filecolor=magenta,
	urlcolor=plotblue,
	citecolor=linkblue,
	pdfpagemode=FullScreen,
}


% Row Height of Table times 1.5
\renewcommand{\arraystretch}{1.5}

% Karriertes Papier für Antwort
\newcommand{\karos}[1]{\begin{tikzpicture}\draw[step=0.5cm,color=gray] (0,0) grid (\linewidth,#1 cm);\end{tikzpicture}}
\graphicspath{{Algebra/2_terme/images/}}

\newcommand{\Skriptname}{Terme und Termumformungen}
\newcommand{\Autor}{Jonas Guler (Gu)}

\begin{document}   %KSH-Gubser - LATEX/PDF
\begin{titlepage}
	\clearpage
	\title{\Huge\textbf{\Skriptname} \\[1cm]}
	\author{Autor:\\ \Autor}
	\date{\normalsize {Version 1.0}\\ \datelong}
	\maketitle
    \thispagestyle{empty}
    \vfill
    \tableofcontents
\end{titlepage}
\pagenumbering{arabic}
%########################################## Kapitel 2 ##################################################
\section{Der Begriff des Terms}
    Wir haben schon eine Vorstellung davon, was ein Term ist. Man kann den Begriff folgendermassen definieren:
    \begin{definition}
        Unter einem Term versteht man eine sinnvolle Aneinanderreihung von endlich vielen mathematischen Symbolen, in der aber kein Relationszeichen vorkommt.\newline
    	Folgende Bestandteile kann ein Term beinhalten:
    	\begin{enumerate}[label=\ ]
    		\item Jede Zahl aus $\N,\Z,\Q,\R$ ist ein Term.
    		\item Jede Variablen $x,y,z,\ldots$ und jeder Parameter $a,b,c,\ldots$ ist ein Term.
    	\end{enumerate}
    	Zwei Terme $S$ und $T$ können durch folgende Operationen zu einem neuen Term verkn\"upft werden:
    	\begin{enumerate}[label=\ ]
    		\item Addition: Der Term $S+T$ wird \emph{Summe} genannt, $S$ heisst \emph{1. Summand}, $T$ heisst \emph{2. Summand}.
    		\item Subtraktion: Der Term $S-T$ wird \emph{Differenz} genannt, $S$ heisst \emph{Minuend}, $T$ heisst \emph{Subtrahend}.
    		\item Multiplikation: Der Term $S\*T$ wird \emph{Produkt} genannt, $S$ heisst \emph{1. Faktor}, $T$ heisst \emph{2. Faktor}.
    		\item Division: Der Term $S:T$ wird \emph{Quotient} genannt, $S$ heisst \emph{Dividend}, $T$ heisst \emph{Divisor}. \\
    		Als Bruch: Der Term $\frac{S}{T}$ wird \emph{Bruch} genannt, $S$ heisst \emph{Zähler}, $T$ heisst \emph{Nenner}.
    		\item Potenzen: Der Term $S^{T}$ wird \emph{Potenz} genannt, $S$ heisst \emph{Basis}, $T$ heisst \emph{Exponent}.
    		\item Wurzeln: Der Term $\sqrt{S}$ wird \emph{(Quadrat-)Wurzel} gennant, $S$ heisst \emph{Radikand}.
    	\end{enumerate}
        \textbf{Wichtig:} $=,<,>,\geq,\leq,\in,\subset,\ldots$ dürfen nicht in einem Term enthalten sein!
    \end{definition}
    \begin{beispiel}
    	Wir beginnen mit den Termen $R=x^{2}, S=3y$. Dann können wir auf verschiedene Weise neue Terme bilden:
    	\begin{align*}
    	   R+S = x^{2} + 3y \qquad \frac{R - 10}{S} = \frac{x^{2} - 10}{3y} 
    	\end{align*}
        Oft gibt man an, welche Variablen in einem Term vorkommen, d.h. \glqq von welchen Variablen ein Term abhängt\grqq. Wir würden beispielsweise schreiben:
        \begin{align*}
        	R(x) = x^{2}, \qquad S(y)=3y, \qquad U(x,y)=xy - 5y
        \end{align*}
    	Terme kann man \emph{auswerten}, indem man für die Variablen Zahlenwerte einsetzt.\newline Für ${T(x)=\sqrt{x^{2}-1}}$ gilt:
        \begin{align*}
        	T(5) &=\sqrt{5^{2} - 1} = \hspace{1.5cm} && T(1) = \hspace{2.5cm} \\[0.1cm]
        	T(7) &= && T(0) = 
        \end{align*}
        Für die Reihenfolge, in welcher diese Operationen ausgeführt werden, gilt:
        \begin{itemize}
            \item Operationen 3. Stufe: Potenzieren, Wurzelziehen.
            \item Operationen 2. Stufe: Multiplikation, Division/Brüche
            \item Operationen 1. Stufe: Addition, Subtraktion
        \end{itemize}
        Gleichwertige Operationen werden von links nach rechts ausgeführt:
        \begin{align*}
            9 - 3 + 8 = (9-3)+8 &= \hspace{3cm}&&&\\
            9 - 3 + 8  \neq 9 - (3+8) &= &&&\\
            9 : 3 \cdot 8 = (9:3) \cdot 8 &= &&&\\
            9 : 3 \cdot 8 \neq 9 : (3 \cdot 8) &= &&&
        \end{align*}
        Um die Reihenfolge der Operationen zu verändern, verwenden wir runde oder eckige Klammern $ ( \quad ), [\quad] $. Insgesamt gilt für die Reihenfolge der Auswertung somit \textbf{Kla}mmern - \textbf{Po}tenzen - \textbf{Pu}nktoperationen - \textbf{Stri}choperationen, kurz \textbf{KlaPoPuStri}.
    \end{beispiel}
    \begin{aufgabe}
        \begin{enumerate}
            \item Welche der folgenden Ausdrücke ist ein Term? Begründe deine Antwort.
            \begin{enumcols}[3]
                \item $2^{10}>1000$
                \item $(a-b+c++$
                \item $a^2-b^2$
                \item $a^2-b^2=(a+b)(a-b)$
                \item $\frac{7}{19}\nin\N$
                \item $1+\frac{x}{1+\frac{x}{1+x}}$
            \end{enumcols}
            \item Lange Zeit wurde der Mythos verbreitet, dass es unmöglich sei, ein Stück Papier mehr als $7-$mal zu falten. Im Jahr $2001$ wiederlegte eine Schülerin diese Mythos, in dem sie einen langen Papierstreifen $11-$mal faltete.\newline
            Sie fand zudem einem Term $T(D,n)$, der zu jeder beliebigen Papierdicke $D\:[\text{in mm}]$ die Mindestlänge des Papierstreifens berechnet, die  nötig ist um den Streifen $n-$mal zu falten: \[T(D,n)=\frac{\pi\*D}{6}\left(2^n+4\right)\left(2^n-1\right)\]
            Welche Mindestlänge muss ein $0.1$mm dicker Papierstreifen haben,wenn er $7-$mal, $10-$mal oder $12-$mal gefaltet werden soll.
        \end{enumerate}
    \end{aufgabe}

\newpage
\section{Gleichheit von Termen}
    Wie wir gesehen haben, kann man für die Variablen in Termen Zahlenwerte einsetzen. Man sagt, man wertet den Term an einer Stelle aus:
    \begin{align*}
    	T(x,y) &= (x-y)^{2} - y^{2} \hspace{8cm} \\
    	T(4,1) &= \\ %(4-1)^{2} - 1^{2} = 9 - 1= 8 \\
    	T(3,5) &= \\ %(3-5)^{2} - 5^{2} = 4 - 24 = -20 \\
    	T(2, 0) &= %(2-0)^{2} - 0^{2} = 4 - 0 = 4
    \end{align*}
    Wir vergleichen nun die beiden Terme $ S(x,y) = x^{2} - 2xy $ und $ T(x,y) = (x-y)^{2} - y^{2} $, indem wir verschiedene Auswertungen machen:\vspace{5pt}\newline
    \begin{minipage}{0.47\linewidth}
        Term: $S(x,y)=x^{2}-2xy$\vspace{5pt}\newline
        \begin{tabular}{m{0.8cm}|m{0.8cm}|m{0.8cm}|m{0.8cm}|m{0.8cm}|p{0.8cm}}
        	$x\backslash y$ & 0 & 1 & 2 & 3 & 4 \\\hline
        	0 & & & & & \\\hline
        	1 & & & & & \\\hline
        	2 & & & & & \\\hline
        	3 & & & & & \\\hline
        	4 & & & & & \\
        \end{tabular}
    \end{minipage}
    \hfill
    \begin{minipage}{0.47\linewidth}
        Term: $T(x,y)=(x-y)^{2}-y^{2}$\vspace{5pt}\newline
        \begin{tabular}{m{0.8cm}|m{0.8cm}|m{0.8cm}|m{0.8cm}|m{0.8cm}|m{0.8cm}}
        	$x\backslash y$ & 0 & 1 & 2 & 3 & 4 \\\hline
        	0 & & & & & \\\hline
        	1 & & & & & \\\hline
        	2 & & & & & \\\hline
        	3 & & & & & \\\hline
        	4 & & & & & \\
        \end{tabular}
    \end{minipage}\vspace{5pt}\newline
    Wir machen folgende Beobachtungen: \vspace{42pt}
    %Für die vorgenommenen Termauswertungen stimmen die beiden Terme überein.
    %Dies ist aber keine Garantie, dass die beiden Terme für alle Werte übereinstimmen.
    \begin{definition}
        Für die (Un-)Gleichheit von zwei Termen gilt.
        \begin{itemize}
        \item Ungleichheit: Stimmen $ S(x) $ und $ T(x) $ für \textbf{eine} Termauswertung nicht überein, so sind die Terme verschieden.
        \item Gleichheit: Stimmen $ S(x) $ und $ T(x) $ für \textbf{jede} Termauswertung überein, so sind die Terme gleich.
        \end{itemize}
    \end{definition}
    Die obige Definition ist zwar geeignet, um zu zeigen, dass zwei Terme verschieden sind:
    \begin{align*}
    	U(x,y) &= (x-y)^{2} \hspace{10cm} \\[0.5cm]
    	V(x,y) &= x^{2} - y^{2}
    \end{align*}
    Offensichtlich ist es aber nicht möglich, \emph{alle Termauswertungen} durchzuführen und so zu zeigen, dass zwei Terme gleich sind. Um Gleichheit zu zeigen, werden wir \emph{Termumformungen} benutzen. Um diese Rechengesetze geht es in den nächsten Abschnitten.
    \newpage


\section{Addition und Subtraktion}
    Wiederholen wir zunächst die Gesetze der Addition:
    \begin{satz}
    	Für die Addition von reellen Zahlen $a,b,c\in\R$ gelten folgende Rechengesetze:
    	\begin{enumerate}
    	   \item Die $0$ ist das Neutralelement der Addition.
    	   \item Die Addition ist kommutativ. 
    	   \item Die Addition ist assoziativ.
    	\end{enumerate}
    \end{satz}
    Wir benutzen diese Gesetze als \textbf{Axiome}. D.h. wir gehen davon aus, dass sie wahr sind und benutzen diese Axiome, um weitere Gesetze herzuleiten. Geometrisch kann man die Gesetze plausibel machen, indem man Zahlen als Pfeile auf der Zahlengerade auffasst und die Addition als das Aneinanderhängen von diesen Pfeilen interpretiert:
    \begin{figure}[H]
        \centering
        \begin{tikzpicture}[line cap=round,line join=round,>=triangle 45,x=1.0cm,y=1cm]
            \draw[->](-5,0) --(8,0);
            \foreach \x in {-4.,-3.,-2.,-1.,0.,1.,2.,3.,4.,5.,6.,7.}
            \draw[shift={(\x,0)},color=black] (0pt,2pt) -- (0pt,-2pt) ;%node[below] {\footnotesize $\x$};
            \draw[color=black] (8,0) node[right] {$\R$};
        \end{tikzpicture}
    \end{figure}\vspace{2cm}

    Die Subtraktion braucht keine eigenen Gesetze, da sie keine neue Operation darstellt. Um das zu erkennen, wiederholen wir den Begriff der \textbf{Gegenzahl}:
    \begin{definition}
    	Die Gegenzahl einer Zahl $a\in\R$ ist diejenige Zahl $x$ mit demselben Betrag und umgekehrten Vorzeichen.\newline
    	Für die Gegenzahl von $a$ schreiben wir: $x=-a$.
    \end{definition}
    \begin{beispiel}
    	Beschreibe formal:
    	\begin{enumerate}[label=\alph*)]
    		\item Das Produkt von zwei Gegenzahlen ist gleich dem Produkt der beiden Zahlen.\vspace{18pt}
    		\item Die Gegenzahl ist gleich dem Produkt aus $ -1 $ und dieser Zahl.\vspace{18pt}
    		\item Die Gegenzahl der Gegenzahl ist wieder gleich der Zahl.\vspace{18pt}
    		\item Die Differenz zweier Zahlen ist gleich der Addition aus Minuend und der Gegenzahl des Subtrahenden.\vspace{18pt}
    		\item Die Gegenzahl einer Summe ist gleich der Summe der Gegenzahlen.\vspace{30pt}
    	\end{enumerate}
    \end{beispiel}
    Punkt d) zeigt insbesondere folgendes zur Subtraktion:
    \begin{satz}
        Die Subtraktion ist keine eigene Operation. Die Subtraktion entspricht der Addition einer Gegenzahl.\vspace{18pt}
    \end{satz}
    Aus den Gesetzen der Addition und der Schreibweise der Subtraktion als Addition der Gegenzahl folgen die Klammerregeln:
    \begin{satz}
    	Für Addition und Subtraktion von reellen Zahlen $a,b,c\in\R$ gelten folgende Klammerregeln:
    	\begin{enumerate}
    		\item $a+(b+c) = a+b+c$ \quad \text{bzw.} \quad $a+(b-c) = a+b-c$
    		\item $a-(b+c) = a-b-c$ \quad \text{bzw.} \quad $a-(b-c) = a-b+c$
    	\end{enumerate}
    \end{satz}
    \begin{aufgabe}
        Löse die Aufgaben $18-25$ aus der Aufgabensammlung.
    \end{aufgabe}

\newpage
\section{Multiplikation und Division}
    Für ganze Zahlen ist die Multiplikation eine Zusammenfassung der Addition:
    \begin{align*}
	   n \* a = \underbrace{a + a + a + \ldots + a}_{n \text{ Summanden}}
    \end{align*}
    Wiederholen wir zunächst die Gesetze der Multiplikation:
    \begin{satz}\label{satz:GesetzeDerMultiplikation}
    	Für die Multiplikation von reellen Zahlen $a,b,c\in\R$ gelten folgende Rechenregeln:
    	\begin{enumerate}
    		\item Die 1 ist das Neutralelement der Multiplikation.
    		\item Die Multiplikation ist kommutativ.
    		\item Die Multiplikation ist assoziativ.
    		\item Die 0 ist der einzige Nullteiler in $ \mathbb{R} $.
    	\end{enumerate}
    \end{satz}
    Wir nehmen diese Gesetze wie schon bei der Addition als gegeben hin. Man kann die Gesetze aber geometrisch veranschaulichen und somit plausibel erklären:
    \begin{figure}[H]
        \centering
        \begin{tikzpicture}[line cap=round,line join=round,>=triangle 45,x=1.0cm,y=1cm]
            \draw[->](-5,0) --(8,0);
            \foreach \x in {-4.,-3.,-2.,-1.,0.,1.,2.,3.,4.,5.,6.,7.}
            \draw[shift={(\x,0)},color=black] (0pt,2pt) -- (0pt,-2pt) ;%node[below] {\footnotesize $\x$};
            \draw[color=black] (8,0) node[right] {$\R$};
        \end{tikzpicture}
    \end{figure}\vspace{1cm}
    Die Division braucht keine eigenen Gesetze, da sie keine neue Operation darstellt. Um das einzusehen, wiederholen wir den Begriff des Kehrwerts:
    \begin{definition}
    	Der Kehrwert einer Zahl $a\in\R$ ist diejenige Zahl $x\in\R$, welche die folgende Gleichung erfüllt:
    	\begin{align*}
    		a\*x=1
    	\end{align*}
    	Für den Kehrwert von $a$ schreiben wir: $x=\frac{1}{a}$. Bemerke, $0$ besitzt keinen Kehrwert!
    \end{definition}
    \begin{beispiel}
    	Schreibe formal:
    	\begin{enumerate}[label=\alph*)]
    		\item Der Kehrwert einer Gegenzahl ist die Gegenzahl des Kehrwerts dieser Zahl. \vspace{18pt}
    		\item Der Kehrwert einer Summe ist \emph{nicht gleich} der Summe der Kehrwerte.\vspace{18pt}
    		\item Der Kehrwert eines Bruches ist gleich dem Bruch dessen Zähler gleich dem ursprünglichen Nenner und dessen Nenner gleich dem ursprünglichen Zähler ist.\vspace{18pt}
    		\item Der Quotient zweier Zahlen ist gleich dem Produkt aus Dividend und dem Kehrwert des Divisors.\vspace{18pt}
    	\end{enumerate}
    \end{beispiel}
    Punkt d) zeigt insbesondere folgendes zur Division:
    \begin{satz}
        Die Division ist keine eigene Operation. Die Division entspricht der Multiplikation mit dem Kehrwert.\vspace{18pt}
    \end{satz}
    Es ist wichtig, dass wir gut verstehen, wie sich die Multiplikation (und damit auch die Division) im Zusammenhang mit Vorzeichen verhalten:
    \begin{align*}
        2\*5 =&\hspace{3cm} &&2\*(-5)=\hspace{5cm} \\
        (-2)\*5=& &&(-2)\*(-5)=\\
        2:5=\frac{2}{5}=& &&2:(-5)=\frac{2}{-5}=\\
        (-2):5=\frac{-2}{5}=& &&(-2):(-5)=\frac{-2}{-5}=
    \end{align*}
    Diese Rechengesetze gelten allgemein:
    \begin{satz}[Punktoperationen und Vorzeichen]\label{satz:PunktVorzeichen}
        \begin{table}[H]
            \centering
            \begin{tabular}{m{0.5cm} m{4.3cm} m{3.2cm} m{0.7cm} m{3.5cm}}
                \multicolumn{5}{l}{Für das Multiplizieren von Termen gelten folgende Regeln:}\\
                1. & Gleiches Vorzeichen & $a\*b = ab$ & bzw. & $(-a)\*(-b) = ab$ \\
                2. & Verschiedene Vorzeichen & $(-a)\*b = -ab$ & bzw. & $a\*(-b) = -ab$ \\
                \multicolumn{5}{l}{Für das Dividieren von Termen unterschiedlicher Vorzeichen gelten analog:}\\
                3. & Gleiches Vorzeichen & $a\div b = \frac{a}{b}$ & bzw. & $(-a)\div(-b) = \frac{-a}{-b}=\frac{a}{b}$ \\
                4. & Verschiedene Vorzeichen & $(-a)\div b = \frac{-a}{b}=-\frac{a}{b}$ & bzw. & $a\div(-b) = \frac{a}{-b}=-\frac{a}{b}$ \\
            \end{tabular}
        \end{table}
    \end{satz}
    \begin{aufgabe}
        \begin{enumerate}
            \item Bestimme folgende Produkte und achte dabei auf das Vorzeichen des Resultats.
            \begin{enumcols}[2]
                \item $(-3)\*2\*(-1)\*(-5)$
                \item $(-3)\*(-b)\*(-2a)\*(-3c)$
                \item $2\*(-1)\*4\*(-2)$
                \item $(-a)\*b^2\*(-ab)\*ab^2$
                \item $5xy\*(2x^2)\*(-2x)^2$
                \item $\frac{1}{3}\*\left(-\frac{6}{5}\right)\*\frac{-7}{4}\*\left(-\frac{10}{3}\right)$
            \end{enumcols}\newpage
            \item Bestimme folgende Potenzen und achte dabei auf das Vorzeichen des Resultats. Welche Gesetzmässigkeiten kannst du feststellen? Stelle eine Regel für Potenzen der Form $(-a)^n$ auf.
            \begin{enumcols}[4]
                \item $(-2)^1$
                \item $(-2)^3$
                \item $(-2)^5$
                \item $(-2)^7$
                \item $(-2)^2$
                \item $(-2)^4$
                \item $(-2)^6$
                \item $(-2)^8$
            \end{enumcols}
            \item Fasse die folgenden Multiplikationen so weit wie möglich zusammen. % Algebra 1 S23 A140
            \begin{enumcols}[2]
                \item $(-5n^2)2n^3(-2n)^3$
                \item $(-7a^3)(-2ab)^3(-a^2b)$
                \item $-5(-12)3a^2b(-4ab^2)^2$
                \item $-\frac{1}{4}\left(-\frac{1}{2}a\right)(-6x)(-y)^2$
            \end{enumcols}
            \item Berechne die folgenden Divisionen und kürze das Resultat so weit wie möglich. % Algerba 1 S24 A156
            \begin{enumcols}[2]
                \item $108a^2\div 9a$
                \item $528a^6bc\div 88ab$
                \item $135c^3\div 15c$
                \item $27ab^2\div 4a$
            \end{enumcols}
        \end{enumerate}
    \end{aufgabe}


\newpage
\section{Das Distributivgesetz}
    Nach dem Gesetzen der Addition und Multiplikation benötigen wir ein Gesetz, welches beide Operationen verbindet: Das Distibutivgesetz (lateinisch distribuere=verteilen).
    \begin{satz}
        Die Multiplikation ist distributiv über der Addition. D.h. für drei rationale Zahlen $a,b,c$ gelten
        \[a\*(b+c)=a\*b+a\*c\]
        \begin{figure}[H]
            \centering
            \begin{tikzpicture}[scale=0.75][line cap=round,line join=round,>=triangle 45,x=1.0cm,y=1.0cm]
                \clip(-3.8,-1.82) rectangle (14.86,2.5);
                \draw [line width=1pt] (-3.,-1.)-- (2.,-1.);
                %\draw [line width=1pt] (2.,-1.)-- (2.,2.);
                \draw [line width=1pt] (2.,2.)-- (-3.,2.);
                \draw [line width=1pt] (-3.,2.)-- (-3.,-1.);
                \draw [line width=1pt] (2.,-1.)-- (4.,-1.);
                \draw [line width=1pt] (4.,-1.)-- (4.,2.);
                \draw [line width=1pt] (4.,2.)-- (2.,2.);
                \draw [line width=1pt] (6.,-1.)-- (11.,-1.);
                \draw [line width=1pt] (11.,-1.)-- (11.,2.);
                \draw [line width=1pt] (11.,2.)-- (6.,2.);
                \draw [line width=1pt] (6.,2.)-- (6.,-1.);
                \draw [line width=1pt] (12.,-1.)-- (14.,-1.);
                \draw [line width=1pt] (14.,-1.)-- (14.,2.);
                \draw [line width=1pt] (14.,2.)-- (12.,2.);
                \draw [line width=1pt] (12.,2.)-- (12.,-1.);
                \draw (-.5,0.86) node[anchor=north west] {$a\*(b+c)$};
                \draw (8.14,0.86) node[anchor=north west] {$a\*b$};
                \draw (12.5,0.86) node[anchor=north west] {$a\*c$};
                \draw (4.68,0.82) node[anchor=north west] {$=$};
                \draw (11.2,0.86) node[anchor=north west] {$+$};
            \end{tikzpicture}
        \end{figure}
    \end{satz}

\subsection{Ausmultiplizieren von Summen}
    Mit dem Distributivgesetz lassen sich auch Produkte von Summen ausmultiplizieren.
    \begin{beispiel}\label{bsp:Distributiv}
    	Entwickle in eine Summe und fasse zusammen:
    	\begin{enumerate}[label=\alph*)]
    		\item $(x-y)(3x-y)=$\vspace{24pt}
    		\item $(-2b-5a)(3a+b)=$\vspace{24pt}
    		\item $(u+2v)(3uv-2v-4u)=$\vspace{24pt}
    		\item $(4a+3s)(-2s+2r-5)=$\vspace{30pt}
    		\item $(2u+v)(v-u)(3u+4v)=$\vspace{30pt}
    	\end{enumerate}
    \end{beispiel}
    \begin{aufgabe}
        Löse von den Aufgaben $27-45$ jeweils alle ungeraden Nummern.
    \end{aufgabe}

\subsection{Die binomischen Formeln}
    Wir repetieren die binomischen Formeln. Wir erklären zunächst den Begriff des Polynoms:
    \begin{definition}
    	Unter dem Begriff \emph{Polynom} (pol\'{y}s griechisch viel, n\'{o}mos griechisch Gesetz/Regel) fasst man eine Familie von Termen zusammen:
    	\begin{align*}
    		&-3uv^{3}w &&\text{Monom oder 1-gliedriges Polynom} \\
    		&2ab - b^{3} &&\text{Binom oder 2-gliedriges Polynom} \\
    		&-x + y - z &&\text{Trinom oder 3-gliedriges Polynom}
    	\end{align*}
    	Auch folgende Terme sind Polynome:
    	\begin{align*}
    		5(x^{2} - 3x +1) = 5x^{2} - 15x + 5 \qquad
    		\frac{a-b}{2} = \frac{1}{2}a - \frac{1}{2} b
    	\end{align*}
    	Folgende Terme sind aber keine Polynome:
    	\begin{align*}
    		2^{x} - 4x, \quad \sqrt{a^{2} + b^{2}}, \quad \frac{1}{a^{2} - 5a + 3}
    	\end{align*}
    \end{definition}
    Aus dem Distributivgesetz folgen nun die drei binomischen Formeln:
    \begin{satz}
	   Für Zahlen $a,b\in\R$ gelten die drei binomischen Formeln:
    	\begin{enumerate}
    		\item binomische Formel: \qquad $(a+b)^{2} = a^{2} + 2ab + b^{2}$
    		\item binomische Formel: \qquad $(a-b)^{2} = a^{2} - 2ab + b^{2}$
    		\item binomische Formel: \qquad $(a-b)(a+b) = a^{2} - b^{2}$
    	\end{enumerate}
    \end{satz}
    \begin{aufgabe}
        Löse von den Aufgaben $47-59$ jeweils alle ungeraden Nummern.
    \end{aufgabe}

\subsection{Höhere Potenzen von $(a+b)$}
    In diesem Abschnitt wollen wir der Frage nachgehen, wie ein Term der Form $(a+b)^n$ berechnet werden kann. Dabei ist der Parameter $n\in\N_0$ eine beliebige natürliche Zahl. Die drei ersten Terme mit $n=1,2,3$ sind uns schon bekannt.
    \begin{align*}
        (a+b)^0=\:&1 && n=0\\
        (a+b)^1=a&+b && n=1\\
        (a+b)^2=a^2+&2ab+b^2 && n=2
    \end{align*}
    \begin{aufgabe}
        \begin{enumerate}
            \item Berechne die untenstehenden Terme, indem du sie ausmultiplizierst.
            \begin{enumerate}[label=\ ]
                \item $(a+b)^3=$
                \item $(a+b)^4=$
                \item $(a+b)^5=$
            \end{enumerate}
            \item Vergleiche nun die Exponenten der beiden Variablen $a$ und $b$ aller $6$ ausmultiplizierter Terme und suche nach einer Gesetzmässigkeit. Was stellst du fest?
        \end{enumerate}
        \karos{2}
        \begin{enumerate}\setcounter{enumi}{1}
            \item Vergleiche nun die Koeffizienten (=Zahl vor den Variablen) aller $6$ ausmultiplizierter Terme. Was fällt die dabei auf? Welcher Zusammenhang besteht zu der untenstehenden Zahlenpyramide?
            \begin{figure}[H]
                \centering
                \includegraphics[scale=0.5]{pascaldreieck.png}
            \end{figure}
            \item Schreibe nun mithilfe der obigen beiden Erkenntnissen die Terme für $(a+b)^6$ und $(a+b)^7$ direkt hin. Ausmultiplizieren ist hier verboten!
            \begin{enumerate}[label=\ ]
                \item $(a+b)^6=$
                \item $(a+b)^7=$
            \end{enumerate}
            \item Multipliziere den untenstehenden Term aus. Was ändert sich dabei? Warum?
            \begin{enumerate}[label=\ ]
                \item $(a\mathbf{-}b)^5=$
            \end{enumerate}
            \karos{2}
        \end{enumerate}
    \end{aufgabe}
    \begin{aufgabe}
        Nutze, die Gesetzmässigkeiten von oben, um auch diese beiden Terme zu notieren. Ausmultiplizieren ist auch hier verboten!
        \begin{enumcols}[2]
            \item $(2x+3y)^3$
            \item $(4r-7s)^6$
        \end{enumcols}
        Löse nun die beiden Aufgaben $73$ und $74$ in der Aufgabensammlung.
    \end{aufgabe}
    

\newpage
\section{Faktorisierungsmethoden}
    Wie wir bereits erwähnt haben, wendet man beim Faktorisieren das Distributivgesetz an, um eine Summe in ein Produkt umzuwandeln. Während ausmultiplizieren \glqq einfaches Handwerk\grqq\: ist, erfordert Faktorisieren etwas mehr Raffinesse. Wir lernen im folgenden vier Faktorisierungsmethoden kennen.

\subsection{Ausklammern von gemeinsamen Faktoren}
    Besitzen alle Summanden einer Summe einen gemeinsamen Faktor, so kann man diesen ausklammern. Der grösstmögliche auszuklammernde Faktor ist dabei gleich dem $\ggT$ der Summanden. Der $\ggT$ bezieht sich hier nicht nur auf Zahlen, sondern auch auf Variabeln.
    \begin{beispiel}
        Klammere den $\ggT$ der Summanden aus:
        \begin{enumerate}[label=\alph*)]
            \item $6rs-18r^{2}=$\vspace{18pt}
            \item $-35xyz-7xz^{2}+14x^{2}yz=$\vspace{18pt}
            \item $27np-12nq+21n=$\vspace{18pt}
            \item $60u^{2}v^{2}-24uv^{2}+48uv-36u^{2}v = $\vspace{18pt}
            \item $-84s^{2}(r+1)+24s^{3}(r+1)-60s(r+1)^{2} = $\vspace{24pt}
        \end{enumerate}
    \end{beispiel}
    \begin{aufgabe}
        Löse von den Aufgaben $75-90$ jeweils die ungeraden Nummern.
    \end{aufgabe}

\subsection{Ausklammern in Teilsummen}
    Die zweite Methode, welche wir studieren, ist das \glqq Ausklammern in Teilsummen\grqq. Im Wesentlichen kehrt man dabei die Schritte des Ausmultiplizierens zweier Summen um.
    \begin{beispiel}
        Im folgenden Beispiel multiplizieren wir ein Produkt von zwei Summen aus. Studiere, wie die vier Summanden entstehen:
        \[(3x-y)(-5a+2b) = 3x(-5a+2b) -y(-5a+2b) = -15ax + 6bx + 5ay - 2by\]
        Faktorisiere nun folgende Summe, indem du die Schritte aus obigen Beispiel in umgekehrter Reihenfolge anwendest:
        \begin{enumerate}[label=\alph*)]
            \item $12ax - 20bx -21ay + 35by = $ \vspace{30pt}
            \item $-18ru - 6rv + 30su + 10sv =$ \vspace{24pt}
            \item $48m^{2}rs + 8rs - 54m^{2}t - 9t = $ \vspace{18pt}
        \end{enumerate}
    \end{beispiel}

    Ein erstes Problem der Methode des Ausklammerns in Teilsummen besteht darin, dass die Faktorisierung nicht eindeutig ist.
    \begin{beispiel}
        Faktorisiere folgende Terme durch Ausklammern in Teilsummen. Durch die Reihenfolge der Summanden werden zwei verschiedene Vorgehen nahegelegt:
        \begin{enumerate}[label=\alph*)]
            \item $18pu + 60qu - 27pv - 90qv = $ \vspace{24pt}
            \item $18pu - 27pv + 60qu - 90qv =$\vspace{18pt}
        \end{enumerate}
        Dieses Problem kann man nun lösen, indem man zunächst den $\ggT$ aller Koeffizienten ausklammert:
        \begin{enumerate}[label=\alph*)]
            \item $-140ax - 56ay + 175bx + 70by = $\vspace{30pt}
            \item $10su + \frac{5}{3}tu - 2sv - \frac{2}{3}tv =$\vspace{18pt}
        \end{enumerate}
    \end{beispiel}
    \begin{aufgabe}
        Löse die Aufgaben $91-96$ aus der Aufgabensammlung.
    \end{aufgabe}

\subsection{Faktorisieren mit den Binomischen Formeln}
    Das stärkste Werkzeug, um eine Summe zu faktorisieren, sind die binomischen Formeln (kombiniert mit Ausklammern). Im Wesentlichen wenden wir dabei die Binomischen Formeln in umgekehrter Richtung an.
    \begin{beispiel} Faktorisiere die folgenden Summen mit Hilfe der binomischen Formeln.
        \begin{enumerate}[label=\alph*)]
            \item $x^{2}-6x+9 =$
            \item $9x^{2}-4y^{2} =$
            \item $16x^{2}+56xy^{2}+49y^{4} =$
            \item $9r^{2}-16s^{6} =$
            \item $9m^{4}+12m^{2}np+4n^{2}p^{2} =$
            \item $25r^{4}s^{2}-9t^{8}u^{2} =$
        \end{enumerate}
    \end{beispiel}
    Bei einigen Summen muss zunächst ein gemeinsamer Faktor ausgeklammert werden, bevor man die binomischen Formeln anwendet.
    \begin{beispiel}
        Faktorisiere die folgenden Summen mithilfe der binomischen Formeln und klammere zuerst einen geeigneten Faktor aus.
        \begin{enumerate}[label=\alph*)]
            \item $75a^{2}-27b^{6} =$
            \item $12u^{4}-75v^{2} =$
            \item $50r^{2}s^{6}-80rs^{3}t^{2}+32t^{4} =$
            \item $-3u^{4}+30u^{2}v-75v^{2} =$
            \item $4m^{3}-16mn^{4} =$
            \item $4x^{2}z+16xyz+16y^{2}z =$
            \item $-125x^{2}y+350xyz^{2}-245yz^{4} =$
        \end{enumerate}
    \end{beispiel}
    \begin{aufgabe}
        Löse die Aufgaben $97-104$ aus der Aufgabensammlung.
    \end{aufgabe}

\subsection{Faktorisieren mit einem Ansatz}
    Polynome in einer Variablen lassen sich mit Hilfe eines \emph{Ansatzes} faktorisieren. Wir wiederholen dazu zunächst den Begriff des Koeffizienten:
    \begin{definition}
	   Ein Polynom der Form
	   \[p(x) = ax^{2} + bx + c \qquad a,b,c\in\R\]
	   nennt man ein \emph{quadratisches Polynom} oder ein \emph{Polynom von Grad 2} in einer Variablen. Die Zahlen $a,b,c$ nennt man \emph{Koeffizienten} des Polynoms. Man sagt, $a$ ist der Koeffizient des quadratischen Terms $ax^{2}$, $b$ ist der Koeffizient des linearen Terms $bx$ und $c$ ist der Koeffizient des konstanten Terms.
    \end{definition}
    Wir studieren nun zunächst wieder das Ausmultiplizieren von zwei Summen:
    \begin{beispiel}
        Beim Ausmultiplizieren entsteht ein quadratischer Term, ein linearer Term und ein konstanter Term. Überlege dir, wie der lineare Koeffizient und der konstante Koeffizient aus den beiden Faktoren entstehen:
        \[(x-7)(x+3) = \hspace{10cm}\]
        Wir halten fest: Der Koeffizient des konstanten Terms ist gleich \gap{5} der Zahlen der beiden Faktoren. Der Koeffizient des linearen Terms ist gleich \gap{4} oder \gap{5} der Zahlen der beiden Faktoren. \newline
        Faktorisiere nun folgende Produkte, indem du den Ansatz $ (x \pm \ldots )(x \pm \ldots) $ machst. Anschliessend zerlegst du den Koeffizienten des konstanten Terms in ein Produkt und überlegst dir, wie du die Faktoren zum Koeffizienten des linearen Terms addieren kannst.
        \begin{enumerate}[label=\alph*)]
            \item $x^{2}-x+12 =$ \vspace{18pt}
            \item $x^{2}-7x+10 =$ \vspace{18pt}
            \item $x^{2}-7x-8 =$ \vspace{18pt}
        \end{enumerate}
        Wir können noch eine weiter Beobachtung festhalten: Hat der Koeffizient des konstanten Terms ein positives Vorzeichen, so muss man den Ansatz $ (x+ \ldots)(x+\ldots)$ oder $(x-\ldots)(x-\ldots)$ verwenden. Hat der Koeffizient des konstanten Terms jedoch ein negatives Vorzeichen, so muss man den Ansatz $(x+\ldots)(x-\ldots)$ verwenden.
    \end{beispiel}
    Ist der Koeffizient des quadratischen Terms nicht mehr =1, so ist das Anwenden der Methode etwas komplexer:
    \begin{beispiel}
        Bei folgenden Polynomen ist der quadratische Koeffizient $ \neq 1 $. Je nach Situation muss die Ansatzmethode anders angepasst werden.
        \begin{enumerate}[label=\alph*)]
            \item $2x^{2} + 4x - 30 = 2(x^{2}+2x-15) =$\vspace{18pt}
            \item $4x^{2} + 40x + 96 = $ \vspace{18pt}
        \end{enumerate}
    \end{beispiel}
    Wir halten also fest: Ist der Koeffizient des quadratischen Terms nicht gleich 1, muss der Koeffizient $a$ ausgeklammert werden können anderen falls sind wir nicht in der Lage den Ausdruck zu faktrieren.
    \begin{aufgabe}
        Löse die Aufgaben $105-111$ aus der Aufgabensammlung.
    \end{aufgabe}
    \begin{aufgabe}[für Schnelle]
        Löse auch die gemischten Aufgaben $112-120$ aus der Aufgabensammlung.
    \end{aufgabe}


\newpage
\subsection*{History}
\begin{table}[H]
    \centering
    \begin{tabular}{|p{13mm}|p{20mm}|p{30mm}|p{70mm}|}\hline
        %\multicolumn{4}{|l|}{History}\\\hline
        Version & Datum & Person & Bemerkung \\\hline
        $0.1$ & 02.08.2024 & Jonas Guler & Kapitel erstellt \\\hline
        $0.2$ & ToDo & Jonas Guler & \\\hline
    \end{tabular}
\end{table}

\end{document}